%!TEX root=../../main.tex

\begin{proposition}\label{prop:min-norm-interp}
    Consider the min group-norm interpolation problem,
    \[
        \min_{w} \sum_{\bi \in \calB} \norm{\wi}_2 \quad \text{s.t.} \, \, X w = y.
    \]
    There exist \( X, y \) such that
    the minimum-group-norm solution to this problem
    is unique and not in \( \Row(\Xas) \).
\end{proposition}
\begin{proof}
    We provide a counter-example where the solution is not the row space of
    the active set.
    Choose
    \[
        X =
        \left[
            \begin{array}{cc|c}
                1 & 2 & 0 \\
                1 & 0 & 2 \\
            \end{array}
            \right],
        \quad
        y = \begin{bmatrix}
            1 \\
            1
        \end{bmatrix},
    \]
    where the vertical line indicates the block structure, i.e.,
    \( b_1 = \cbr{1, 2} \) and \( b_2 = \cbr{3} \).

    Clearly a solution using only \( b_2 \) cannot interpolate the data,
    so the active set must be \( \cbr{b_1, b_2} \)  or \( \cbr{b_1} \).
    If the active set is \( b_1 \), then the minimum norm interpolating solution
    can only be \( w = [1 \; 0 \; 0]^\top \), which has group norm \( 1 \).

    Now, consider when the active set is \( \cbr{b_1, b_2} \).
    The interpolating solution in \( \Row(X) \) satisfies the following
    system
    \begin{align*}
        \left[
            \begin{array}{cc|c}
                1 & 2 & 0 \\
                1 & 0 & 2 \\
            \end{array}
            \right]
        \rbr{
            \alpha
            \begin{bmatrix}
                1 \\ 2\\ 0
            \end{bmatrix}
            +
            \beta
            \begin{bmatrix}
                1 \\ 0\\ 2
            \end{bmatrix}
        }
         & = \begin{bmatrix}
                 1 \\
                 1
             \end{bmatrix} \\
        \implies
        \begin{bmatrix}
            5 \alpha + \beta \\
            \alpha + 5 \beta
        \end{bmatrix}
         & =
        \begin{bmatrix}
            1 \\
            1
        \end{bmatrix}.
    \end{align*}
    Solving for \( \alpha \) and \( \beta \) yields
    \( \alpha = 1 - 5 \beta \) and \( 24 \beta = 4 \),
    which implies \( \beta = 1/6 \) and \( \alpha = 1/6 \).
    The optimal \( w^* \) is thus
    \[
        w^* =
        \begin{bmatrix}
            1/3 \\
            1/3 \\
            1/3
        \end{bmatrix},
    \]
    and the group norm of \( w^* \) is
    \[
        \sum_{\bi \in \calB} \norm{\wi^*}_2 = \sqrt{2/9} + 1/3  = (1 + \sqrt{2}) / 3 < 1.
    \]
    Thus, we can discount the interpolating solution using only \( b_1 \).

    Now let's see if we can reduce the norm by including directions in \( \Null(X) \).
    The Null space is orthogonal to both rows of \( X \), from which we conclude
    \[
        \Null(X) =
        \cbr{
            \gamma z : z =
            \begin{bmatrix}
                -2/3 \\
                1/3  \\
                1/3
            \end{bmatrix},
            \gamma \in \R
        }.
    \]
    Any vector \( w' = w^* + \gamma z \) is an interpolating solution,
    so it only remains to check if there is a choice of \( \gamma \)
    that decreases the group norm.
    Assuming \( \gamma > -1 \),
    \begin{align*}
        \sum_{\bi \in \calB} \norm{\wi^*}_2
         & \leq \sum_{\bi \in \calB} \norm{\wi'}_2                                                                             \\
        \iff \frac{1 + \sqrt{2}}{3}
         & = \sqrt{(\frac{1}{3} - \frac{2 \gamma}{3})^2 + (\frac{1}{3} + \frac{\gamma}{3})^2}
        + \abs{\frac{1 + \gamma}{3}}                                                                                           \\
        \iff \frac{1 + \sqrt{2}}{3} - \frac{1 + \gamma}{3}
         & \leq \sqrt{(\frac{1}{3} - \frac{2 \gamma}{3})^2 + (\frac{1}{3} + \frac{\gamma}{3})^2} \tag{since \( \gamma > -1 \)} \\
        \iff \rbr{\frac{1 + \sqrt{2}}{3} - \frac{1 + \gamma}{3}}^2
         & \leq (\frac{1}{3} - \frac{2 \gamma}{3})^2 + (\frac{1}{3} + \frac{\gamma}{3})^2                                      \\
        \iff \rbr{\sqrt{2} - \gamma}^2
         & \leq (1 - 2 \gamma)^2 + (1 + \gamma)^2
    \end{align*}
    The left-hand side satisfies
    \[
        \rbr{\sqrt{2} - \gamma}^2 = 2 - 2\sqrt{2}\gamma + \gamma^2,
    \]
    while the right-hand side is
    \[
        (1 - 2 \gamma)^2 + (1 + \gamma)^2 = 1 - 4 \gamma + 4 \gamma^2 + 1 + 2 \gamma + \gamma^2
        =  2 - 2 \gamma + 5 \gamma^2.
    \]
    As a result,
    \begin{align*}
        \sum_{\bi \in \calB} \norm{\wi^*}_2
                                                                          & \leq \sum_{\bi \in \calB} \norm{\wi'}_2 \\
        \iff 2 - 2 \gamma + 5 \gamma^2 - 2 - \gamma^2 + 2 \sqrt{2} \gamma & \geq 0                                  \\
        \iff 2 (\sqrt{2} - 1)\gamma + 4 \gamma^2                          & \geq 0.
    \end{align*}
    However, it is easy to check that this fails for \( \gamma \in ((1 - \sqrt{2})/2, 0) \).
    So the minimum-group-norm interpolating solution is not in \( \Row(\Xas) \).

    Now let's show that the minimum-group-norm solution is unique.
    This holds if the group norm as a function of \( \gamma \),
    \[
        g(\gamma) = \sqrt{(\frac{1}{3} - \frac{2 \gamma}{3})^2 + (\frac{1}{3} + \frac{\gamma}{3})^2}
        + \abs{\frac{1 + \gamma}{3}},
    \]
    is strictly convex.
    The function \( h_2(\gamma) = \abs{\frac{1 + \gamma}{3}} \) is convex,
    so we need only show that
    \[
        h_1(\gamma) = \sqrt{(\frac{1}{3} - \frac{2 \gamma}{3})^2 + (\frac{1}{3} + \frac{\gamma}{3})^2}
        = \frac{1}{3} \sqrt{2 - 2\gamma + 5 \gamma^2},
    \]
    is strictly convex.
    Taking the second derivative, we find
    \[
        h_1''(\gamma) = \frac{3}{\rbr{2 - 2\gamma + 5 \gamma^2}^{3/2}} > 0
    \]
    for all \( \gamma \).
    This completes the proof.
\end{proof}

\begin{lemma}\label{lemma:lambda-convergence}
    Let \( W_{g} \) be the set of least squares solutions with minimum group norm.
    That is,
    \[
        W_{g} = \argmin_{w} \cbr{\sum_{\bi \in \calB} \norm{\wi}_2 : X^\top X w = X^\top y }.
    \]
    Then every limit point of the min-norm group lasso solution
    as \( \lambda \rightarrow 0 \) lies in \( W_g \).
\end{lemma}
\begin{proof}
    Let \( \lambda_k \rightarrow 0 \)
    and observe that \( \wmin(\lambda_k) \) has at least one limit point
    since it is bounded (\cref{lemma:bounded-solutions-cgl}).
    Since \( \norm{\ci(\lambda_k)}_2 \leq \lambda_k \),
    we see that \( \lim_{k} \norm{\ci(\lambda_k)}_2 = 0 \)
    and thus \( \lim_k \ci(\lambda_k) = 0 \).

    FO optimality conditions imply
    \begin{equation}
        (X^\top X) \wmin(\lambda) = X^\top y - c(\lambda),
    \end{equation}
    which, taking limits on both sides,
    gives
    \begin{equation}
        \lim_k (X^\top X) \wmin(\lambda_k) = X^\top y.
    \end{equation}
    That is, every limit point \( \bar w \) of \( \wmin(\lambda_k) \) is
    a least squares solution satisfying \( h(\bar w) \leq h(w_g) \).
    We conclude that \( \bar w \in W_g \).
\end{proof}

\rowCE*
\begin{proof}
    Consider the setting of \cref{prop:min-norm-interp},
    with
    \[
        X =
        \left[
            \begin{array}{cc|c}
                1 & 2 & 0 \\
                1 & 0 & 2 \\
            \end{array}
            \right],
        \quad
        y = \begin{bmatrix}
            1 \\
            1
        \end{bmatrix},
    \]
    where the vertical line indicates the block structure, i.e.,
    \( b_1 = \cbr{1, 2} \) and \( b_2 = \cbr{3} \).
    We have shown that the min group norm interpolant is unique,
    is supported on \( b_1 \) and \( b_2 \), and does not lie
    in \( Row(X) \).
    Let \( w_g \) be this solution.

    Let \( \lambda_k \downarrow 0 \).
    By \cref{lemma:lambda-convergence}, every limit point of
    \( \wmin(\lambda_k) = w_g \).
    Thus, \( \lim_k \wmin(\lambda_k) \) exists and is exactly
    \( w_g \).
    Moreover, \( \wmin(\lambda_k) \) must be supported
    on \( b_1 \) and \( b_2 \) for all \( k \) sufficiently large.

    Decomposing \( w_g = a + b \) and \( \wmin(\lambda_k) = r_k + n_k \)
    where \( a, r_k \in \Row(X) \) and \( b, n_k \in \Null(X) \), we see
    that
    \[
        \norm{w_g - \wmin(\lambda_k)}_2^2
        = \norm{a - r_k}_2^2
        + \norm{b - n_k}_2^2 \rightarrow 0,
    \]
    implying that \( n_k \neq 0 \) for sufficiently large \( k \).
    In other words, the min-norm solution to the group lasso problem
    fails to fall in \( \Row(X) \) for some \( \lambda > 0 \).
\end{proof}

\minNormProgram*
\begin{proof}
    Let \( w \in \solfn(\lambda) \).
    By \cref{prop:sol-fn-cgl}, \( \wi = \alpha_\bi \vi \)
    where \( \alpha_\bi \geq 0 \).
    Moreover, \( \alpha_\bi = 0 \) for every
    \( \bi \in \calB \setminus \calS_\lambda \).
    As a result,
    \begin{align*}
        \norm{\w}^2_2
         & = \norm{\we}^2_2                                         \\
         & = \sum_{\bi \in \calS_\lambda} \norm{\alpha_\bi \vi}^2_2 \\
         & = \lambda \norm{\alpha}_2^2,
    \end{align*}
    where the last equality follows from \( \bi \in \calS_\lambda \implies \bi \in \equi \),
    which enforces \( \norm{\vi}_2 = \lambda \).

    Now suppose \( \alpha^* \) is optimal for the cone program
    and let \( w \in \R^d \) such that \( \wi = \alpha_\bi^* \vi \).
    By construction, \( \alpha^*_\bi = 0 \)
    for all \( \bi \in \calB \setminus \calS_\lambda \)
    (or it could not be optimal)
    so that \( \wi = 0 \) for all \( \bi \in \calB \setminus \calS_\lambda \).
    Moreover,
    \[
        X w = \sum_{\bi \in \calS_\lambda} \alpha_\bi^* \Xbi \vi = \hat y,
    \]
    Thus, \( w \) solves
    \[
        \begin{aligned}
            \argmin_{w} \norm{w}_2^2
            \, \, \text{ s.t.} \, \,
             & \forall \, \bi \in \calS_\lambda, \, \wi = \alpha_\bi \vi, \alpha_\bi \geq 0, \, \\
             & \forall \, b_j \in \calB \setminus \calS_\lambda, \w_{b_j} = 0, \,
            X w = \hat y.
        \end{aligned}
    \]
    Invoking \cref{prop:sol-fn-cgl} now proves that \( w \) is the min-norm solution.
\end{proof}

\begin{restatable}{lemma}{l2Reformulation}\label{lemma:l2-reformulation}
    The \( \ell_2 \)-penalized group lasso problem in \cref{eq:l2-penalized-cgl}
    is equivalent to the following CGL problem:
    \begin{equation}\label{eq:l2-reformulation}
        \begin{aligned}
            \min_{w} \,
             & \half \norm{\tilde X w - \tilde y}_2^2 + \lambda \sum_{\bi \in \calB} \norm{\wi}_2 \\
             & \quad \text{s.t.} \quad \Ki^\top \wi \leq 0  \text{ for all } \bi \in \calB.
        \end{aligned}
    \end{equation}
    where we have defined the extended data matrix and targets
    \[
        \tilde X =
        \begin{bmatrix}
            X \\
            \sqrt{\delta} I
        \end{bmatrix}
        \quad \quad
        \tilde y =
        \begin{bmatrix}
            y \\
            0
        \end{bmatrix}.
    \]
    Moreover, \( \tilde X \) is full column-rank and thus the CGL solution is
    unique.
\end{restatable}
\begin{proof}
    It is straightforward to show the equivalence by direct calculation.
    For any \( w \in \R^d \),
    \begin{align*}
        \half \norm{\tilde X \w - \tilde y}_2^2
         & = \half \norm{X w - y}_2^2 + \half \norm{\sqrt{\delta} I w - 0}_2^2 \\
         & = \half \norm{X w - y}_2^2 + \frac{\delta}{2} \norm{w}_2^2,
    \end{align*}
    Substituting this identity into \cref{eq:l2-reformulation} establishes the
    equivalence.

    It is clear by inspection that \( \tilde X \) is full column rank.
    Then \( \Null(X_\calC) = \emptyset \) for all \( \calC \subset \calB \)
    and the solution is unique by \cref{prop:sol-fn-cgl}.
\end{proof}

\penConvergence*
\begin{proof}
    First we show that \( \norm{w^\delta(\lambda)}_2 \leq \norm{w^*(\lambda)}_2 \).
    Suppose by way of contradiction that
    \( \norm{w^\delta(\lambda)}_2 > \norm{w^*(\lambda)}_2 \)
    for some \( \delta > 0 \).
    Since
    \begin{align*}
        \half \norm{Xw^*(\lambda) - y}_2^2
        + \lambda \sum_{\bi \in \calB} \norm{\wi^*(\lambda)}_2
         & =
        \min_{w : \Ki \wi \leq 0} \half \norm{Xw - y}_2^2
        + \lambda \sum_{\bi \in \calB} \norm{\wi}_2 \\
         & \leq
        \half \norm{Xw^\delta(\lambda) - y}_2^2
        + \lambda \sum_{\bi \in \calB} \norm{\wi^\delta(\lambda)}_2,
    \end{align*}
    we deduce
    \begin{align*}
         & \half \norm{Xw^*(\lambda) - y}_2^2
        + \lambda \sum_{\bi \in \calB} \norm{\wi^*(\lambda)}_2
        + \frac{\delta}{2} \norm{w^*(\lambda)}_2^2 \\
         & \hspace{4em} <
        \half \norm{Xw^\delta(\lambda) - y}_2^2
        + \lambda \sum_{\bi \in \calB} \norm{[w^\delta(\lambda)]_\bi}_2
        + \frac{\delta}{2} \norm{w^\delta(\lambda)}_2^2,
    \end{align*}
    which contradicts optimality of \( w^\delta(\lambda) \).
    So \( \norm{w^\delta(\lambda)}_2 \leq \norm{w^*(\lambda)}_2 \)
    for all \( \delta > 0 \).
    As a result, the sequence
    \( \cbr{w^{\delta_k}(\lambda)}_{\delta_k} \), where \( \delta_k \downarrow 0 \),
    is bounded and admits at least one convergent subsequence.
    Let \( \bar w(\lambda) \) be the limit point associated with one
    such subsequence;
    clearly \( \norm{\bar w(\lambda)}_2 \leq \norm{w(\lambda)^*}_2 \).

    Let's show that \( \bar w(\lambda) \) is a solution to the group lasso
    problem by checking the KKT conditions.
    Suppose \( \lambda > 0 \).
    Stationarity of the Lagrangian is
    \[
        X^\top (Xw^{\delta_k}(\lambda) - y) + K \rho^{\delta_k}(\lambda) + s^{\delta_k}(\lambda) + \delta_k w^{\delta_k}(\lambda) = 0,
    \]
    where \( s_\bi^{\delta_k}(\lambda) \in \partial \lambda \norm{\wi^{\delta_k}}_2 \).
    Since \( \norm{s_\bi^{\delta_k}(\lambda)}_2 \leq \lambda \) and \( w^{\delta_k}(\lambda) \)
    is bounded, clearly \( K \rho^{\delta_k}(\lambda) \) is also bounded.

    Dropping to a subsequence if necessary, let \( \lim_{k} w^{\delta_k}(\lambda) = \bar w \)
    and \( \lim_k \Ki \ri^{\delta_k} = \bar{z}_{\bi} \).
    Define
    \[
        R_{1/n} = \cbr{\ri : \norm{\Ki \ri - \bar{z}_{\bi}}_\infty \leq \frac{1}{n}, \ri \geq 0}.
    \]
    The sequence of sets \( R_{1/n} \) is polyhedral and thus retractive.
    Moreover, for each \( n \in \bbN \), there exists \( k \) such that
    \[ \norm{\Ki \ri^{\delta_k} - \bar{z}_{\bi}}_\infty \leq 1 / n, \]
    since \( \Ki \ri^{\delta_k} \rightarrow \bar{z}_{\bi} \).
    Recalling \( \ri^{\delta_k} \geq 0 \) shows that \( \R_{1/n} \) is non-empty.
    The limit of a sequence of nested, non-empty, retractive sets is also
    non-empty \citep[Proposition 1.4.10]{bertsekas2009convex}.
    Moreover, since the limit is exactly
    \[
        \bar R = \cbr{\ri : \Ki \ri = \bar{z}_{\bi}, \ri \geq 0},
    \]
    we deduce that there exists \( \bar \rho \geq 0 \) such that
    \( \Ki \bar \ri = \bar{z}_{\bi} \).

    Taking limits on either side of the stationarity condition, we find
    \begin{align*}
        \Xbi^\top (y - X\bar w(\lambda)) - \Ki \bar \rho(\lambda) = \bar s_\bi,
    \end{align*}
    where the limit point \( \bar s_\bi \) satisfies \( \norm{\bar s_\bi} \leq \lambda \).
    If \( \bi \in \calB \setminus \act(\bar w(\lambda)) \),
    then \( \bwi \) satisfies stationarity.

    Let \( \bi \in \act(\bar w(\lambda)) \).
    Since \( \w^{\delta_k}(\lambda) \rightarrow \bar w(\lambda) \),
    \( \norm{\wi^{\delta_k} - \bar \wi}_2 \rightarrow 0 \)
    and it must happen that
    \( \wi^{\delta_k} > 0 \) for all \( k \) sufficiently large.
    That is,
    \( \calA(w^{\delta_k}(\lambda)) \supseteq \calA(\bar w) \)
    for all \( k \geq k' \).
    Using \( \bi \in \calA(w^{\delta_k}(\lambda)) \)
    provides a closed-form expression for \( s_\bi^{\delta_k} \):
    \begin{align*}
        \lim_{k} s^{\delta_k}(\lambda)
         & = \lambda \lim_{k} \frac{\wi^{\delta_k}(\lambda)}{\norm{\wi^{\delta_k}(\lambda)}_2} \\
         & = \lambda \frac{\bwi(\lambda)}{\norm{\bar \wi(\lambda)}_2},
    \end{align*}
    which shows that \( \bar s_\bi \) is a subgradient of \( \lambda \norm{\wi}_2 \).
    We conclude that the Lagrangian is stationary in \( \bwi \) as well.

    Let us check the remainder of the KKT conditions.
    For feasibility, it is straightforward to observe that
    \[
        \Ki^\top \wi^{\delta_k}(\lambda) \leq 0 \quad \forall \, k \implies \Ki^\top \bwi(\lambda) \leq 0.
    \]
    Similarly,
    \[
        \abr{\ri^{\delta_k}, \Ki^\top \wi^{\delta_k}} = 0 \quad \forall \, k \implies \abr{\bar \ri, \Ki^\top \bwi} = 0,
    \]
    which, combined with \( \bar \rho \geq 0 \), is sufficient to establish
    complementary slackness.
    We have shown the subsequential limits \( (\bar w, \bar \rho) \)
    satisfies the KKT conditions and thus \( \bar w \)
    is a solution to the constrained group lasso problem.

    Since the min-norm solution is unique and
    \( \norm{\bar w(\lambda)}_2 \leq \norm{w^*(\lambda)}_2 \),
    it must be that \( \bar w(\lambda) = w^*(\lambda) \).
    Noting that this holds for every limit point implies
    \( \lim_{\delta \downarrow 0} w^{\delta}(\lambda) \) exists and is \( w^*(\lambda) \).
    This completes the proof for \( \lambda > 0 \).

    If \( \lambda = 0 \), then the proposition follows similarly with
    the additional observation that \( s^{\delta_k}(0) = 0 \)
    for all \( k \).
\end{proof}
